\section{Experiments}
\label{sec:experiments}

In our experiments, we use Gaussian process priors for $f$ with squared-exponential
kernels $k(\x,\x')=\gamma^2\exp\{-0.5 \sum_i (x_i-x_i')^2/\ell_i^2\}$.
The corresponding spectral density is zero-mean
Gaussian with covariance given by $\diag([\ell_i^{-2}])$ and normalizing constant $\alpha=\gamma^2$.
The model hyperparameters are $\{\gamma,\ell_1,\dots,\ell_d,\sigma^2\}$. 
We use broad, uninformative Gamma hyperpriors. 

First, we analyze the accuracy of PES in the task of approximating the differential entropy
(\ref{eq:originalObjective}).  We compare the PES approximation
(\ref{eq:marginalizedObjective}),
with the approximation used by the entropy
search (ES) method \cite{HennigSchuler:2012}.  We also compare with the ground
truth for (\ref{eq:originalObjective}) obtained using a rejection sampling (RS)
algorithm based on (\ref{eq:newObjective}).  For this experiment we generate the
data $\D_n$ using an objective function $f$ sampled from the Gaussian
process prior. The domain $\X$ of $f$ is fixed to be $[0,1]^2$
and data are generated using $\gamma^2=1$, $\sigma^2=10^{-6}$, and
$\ell_i^2=0.1$. To compute (\ref{eq:marginalizedObjective}) we avoid
sampling the hyperparameters and use the known values directly.
We further fix $M=200$ and $m = 1000$.

The ground truth rejection sampling scheme works as follows.  First, $\X$ is discretized
using a uniform grid. The expectation with respect to $p(\xopt|\D_{n})$ in
(\ref{eq:newObjective}) is then approximated using sampling. For this, we sample
$\xopt$ by evaluating a random sample from $p(f|\D_n)$ on each grid cell and
then selecting the cell with highest value.  Given $\xopt$, we then approximate
$\H[p(y|\D_{n},\x,\xopt)]$ by rejection sampling. We draw samples from
$p(f|\D_n)$ and reject those whose corresponding grid cell with highest value is
not $\xopt$.  Finally, we approximate $\H[p(y|\D_{n},\x,\xopt)]$ by first,
adding zero-mean Gaussian noise with variance $\sigma^2$ to the the evaluations
at $\x$ of the functions not rejected during the previous step and
second, we estimate the differential entropy of the resulting samples using
kernels \cite{Ahmad:1976}.

Figure~\ref{fig:costFunctions} shows the objective functions produced by RS, ES
and PES for a particular $\D_n$ with 10 measurements whose locations are
selected uniformly at random in $[0,1]^2$.  The locations of the collected
measurements are displayed with an ``x'' in the plots.  The particular objective function
used to generate the measurements in $\D_n$ is displayed in the left part of
Figure \ref{fig:resultsSyntheticFunctions}.  The plots in Figure
\ref{fig:costFunctions} show that the PES approximation to
(\ref{eq:originalObjective}) is more similar to the ground truth given by RS
than the approximation produced by ES.  In this figure we also see a
discrepancy between RS and PES at locations near $\x= (0.572, 0.687)$. This
difference is an artifact of the discretization used in RS. By
zooming in and drawing many more samples we would see the same behavior in both
plots.

\begin{figure}
\begin{center}
\includegraphics[width=0.9\linewidth]{figures/testAccuracyApproximations/figure.pdf}
\end{center}
\vspace{-0.5em}
\caption{\small Comparison of different estimates of the objective function $\alpha_n(\x)$
given by (\ref{eq:originalObjective}). Left, ground truth
obtained by the rejection sampling method RS.
Middle, approximation produced by the ES method. Right, approximation produced
by the proposed PES method. These plots show that the PES objective is much more
similar to the RS ground truth than the ES objective.}\label{fig:costFunctions}
\vspace{-0.5em}
\end{figure}
